\chapter{A Brief Introduction to Mathilda}
\label{ch:intro}

This chapter is a whirlwind tour. Rather than explain any one thing thoroughly---%
that is the job of Chapter~\ref{ch:mathematics} and beyond---it moves quickly across
the landscape, showing a small, concrete example of each of the things Mathilda can
do: arithmetic and numerics, number theory, algebra, calculus, numerical calculus,
linear algebra, and programming. The aim is to give you a feel for the system and a sense of how much
ground it covers, so that the detailed chapters later feel like a return to places
you have already visited.

Everything below is real: each \texttt{In[\,]} was fed to the Mathilda binary\index{In[] and Out[]} as this
book was built and each \texttt{Out[\,]} is exactly what came back. Type them yourself
and you will see the same results. Throughout, each input and its result stand on the
left, with a few words about that line alongside on the right. This chapter is a
starting point and will grow over time; think of it as the first, shortest pass over
material the rest of the book returns to in depth.

\section{Numbers, exact and approximate}
\label{sec:intro-numbers}

Mathilda distinguishes sharply between \emph{exact} numbers\index{number!exact vs approximate} and \emph{approximate}
ones, and it knows how they combine. Watch what happens as we add numbers of different
kinds:

\begin{codepairs}
\pair{03-introduction/numbers/1}{Two integers add to an integer.}
\pair{03-introduction/numbers/2}{An integer plus a rational\index{number!rational} stays \emph{exact}: the
fraction $7/3$, not a decimal.}
\pair{03-introduction/numbers/3}{A rational plus a \emph{real} gives a real\index{number!real}. One
approximate number makes the whole result approximate---a rule called
\emph{contagion}\index{numeric contagion}.}
\pair{03-introduction/numbers/4}{Integers grow without bound; $2^{64}$ is kept
exactly, far beyond what a single machine word can hold.}
\pair{03-introduction/numbers/5}{To get a decimal you ask, with \B{N}: here a quick
machine-precision value of $\sqrt2$, six figures by default.}
\pair{03-introduction/numbers/6}{Ask \B{N} for more digits and it obliges: here is
$\sqrt2$ to thirty places.}
\pair{03-introduction/numbers/7}{And here is $\pi$ to fifty digits, computed to
arbitrary precision.}
\end{codepairs}

Section~\ref{sec:arithmetic} takes up how these three worlds---machine numbers, exact
integers and rationals, and arbitrary-precision reals---fit together.

\section{Elementary functions}
\label{sec:intro-elementary}

Mathilda knows the elementary functions\index{elementary functions}---powers and roots, the exponential and
logarithm, and the trigonometric functions and their inverses---and treats them the
same way it treats numbers: exactly when it can, symbolically when it must, and
numerically when you ask. At arguments where a simple closed form exists, it returns
that exact value:

\begin{codepairs}
\pair{03-introduction/elementary-exact/1}{A perfect square gives an integer.}
\pair{03-introduction/elementary-exact/2}{\B{Sin} of $\pi/6$ is the rational $1/2$.}
\pair{03-introduction/elementary-exact/3}{\B{Cos} of $\pi/4$ comes back as an exact
surd.}
\pair{03-introduction/elementary-exact/4}{So does \B{Tan} of $\pi/3$.}
\pair{03-introduction/elementary-exact/5}{\B{Log} and \B{Exp} undo each other, so
$\log e^{5}=5$.}
\pair{03-introduction/elementary-exact/6}{\B{ArcTan} of $1$ is the exact angle
$\pi/4$.}
\pair{03-introduction/elementary-exact/7}{At an exact argument with \emph{no} closed\index{closed form}
form, \B{Sin}$[1]$ is held as \texttt{Sin[1]}---not silently turned into a decimal.}
\end{codepairs}

When you do want a number, you get one either by supplying an approximate argument or
by applying \B{N}. By default the result is \emph{machine precision}\index{machine precision}---the fast,
fixed-size floating point of the hardware---but \B{N} will also work to
\emph{arbitrary} precision\index{precision!arbitrary}:

\begin{codepairs}
\pair{03-introduction/elementary-numeric/1}{An approximate argument \texttt{1.0}
forces a machine-precision result.}
\pair{03-introduction/elementary-numeric/2}{Applying \B{N} to the exact
\texttt{Sin[1]} gives the same value.}
\pair{03-introduction/elementary-numeric/3}{Ask for forty digits and you get $\sin 1$
to that precision.}
\pair{03-introduction/elementary-numeric/4}{Here is the constant $e$ to forty digits.}
\pair{03-introduction/elementary-numeric/5}{And here is $\ln 2$ to forty digits. The
elementary functions return in Section~\ref{sec:calculus}.}
\end{codepairs}

\section{Special functions}
\label{sec:intro-special}

Beyond the elementary functions lies a second tier---the \emph{special functions}\index{special functions} of
analysis and mathematical physics: the gamma and zeta functions, the error function,
the Bessel and Airy functions, and many more. Mathilda handles them with exactly the
same discipline, returning exact values wherever they exist:

\begin{codepairs}
\pair{03-introduction/special-exact/1}{\B{Gamma} extends the factorial: $\Gamma(5)=4!
=24$.}
\pair{03-introduction/special-exact/2}{At the half-integers it takes the celebrated
value $\Gamma(1/2)=\sqrt\pi$.}
\pair{03-introduction/special-exact/3}{\B{Zeta} gives Euler's $\zeta(2)=\pi^{2}/6$.}
\pair{03-introduction/special-exact/4}{By analytic continuation\index{analytic continuation}, the notorious
$\zeta(-1)=-\tfrac{1}{12}$.}
\pair{03-introduction/special-exact/5}{\B{Beta} reduces to a rational.}
\pair{03-introduction/special-exact/6}{\B{PolyGamma}, the digamma at $1$, is exactly
$-\gamma$, the negative Euler--Mascheroni constant\index{Euler--Mascheroni constant} \B{EulerGamma}.}
\pair{03-introduction/special-exact/7}{Even \B{AiryAi} at the origin has a closed
form, in terms of \B{Gamma}.}
\end{codepairs}

Where no closed form exists, these functions are computed numerically, to machine or
arbitrary precision just like the elementary ones:

\begin{codepairs}
\pair{03-introduction/special-numeric/1}{$\zeta(3)$, Ap\'ery's constant, to machine
precision.}
\pair{03-introduction/special-numeric/2}{The same constant to forty digits.}
\pair{03-introduction/special-numeric/3}{\B{Erf} at an approximate argument.}
\pair{03-introduction/special-numeric/4}{\B{BesselJ} likewise.}
\pair{03-introduction/special-numeric/5}{$\Gamma(1/2)$ to thirty digits, agreeing with
$\sqrt\pi$. Section~\ref{sec:specialfns} is devoted to these functions.}
\end{codepairs}

\section{Number theory}
\label{sec:intro-numbertheory}

Whole numbers have a rich structure, and Mathilda knows a great deal of it. Primes,
factorizations, divisors, greatest common divisors\index{greatest common divisor}, and modular arithmetic\index{modular arithmetic} are all a
single call away.

\begin{codepairs}
\pair{03-introduction/number-theory/1}{\B{Prime} gives the hundredth prime.}
\pair{03-introduction/number-theory/2}{\B{FactorInteger} breaks $360$ into prime
powers $2^{3}\cdot 3^{2}\cdot 5$.}
\pair{03-introduction/number-theory/3}{\B{Divisors} lists every divisor of $36$.}
\pair{03-introduction/number-theory/4}{\B{GCD} is the greatest common divisor.}
\pair{03-introduction/number-theory/5}{\B{ExtendedGCD} adds B\'ezout's coefficients\index{Bezout's coefficients@B\'ezout's coefficients}:
$12 = 1\cdot 48 + (-1)\cdot 36$.}
\pair{03-introduction/number-theory/6}{\B{PowerMod} finds $7^{100}\bmod 13$ without
ever forming $7^{100}$.}
\pair{03-introduction/number-theory/7}{\B{LCM} is the least common multiple;
$\mathrm{lcm}(12,18,20)=180$.}
\pair{03-introduction/number-theory/8}{\B{EulerPhi} counts the integers from $1$ to $36$
that are coprime to it: $\varphi(36)=12$.}
\pair{03-introduction/number-theory/9}{\B{MultiplicativeOrder} gives the least $k$ with
$2^{k}\equiv 1\pmod{11}$: here that order is $10$.}
\pair{03-introduction/number-theory/10}{\B{PrimitiveRoot} returns a generator of the
units modulo~$11$; the least is $2$, which is exactly why its order above is the full
$\varphi(11)=10$.}
\pair{03-introduction/number-theory/11}{\B{Fibonacci} returns the fiftieth Fibonacci
number exactly. Section~\ref{sec:numbertheory} takes up this material.}
\end{codepairs}

\section{Algebra}
\label{sec:intro-algebra}

Algebra is where a computer algebra system\index{computer algebra system} earns its name: it manipulates
expressions containing symbols\index{expression!symbolic}, not just numbers. Two everyday operations are
multiplying out and its partial inverse, collecting like terms\index{collecting like terms}.

\begin{codepairs}
\pair{03-introduction/algebra-expand/1}{\B{Expand} multiplies out a power.}
\pair{03-introduction/algebra-expand/2}{It works in several variables too.}
\pair{03-introduction/algebra-expand/3}{\B{Collect} gathers terms by the powers of a
chosen variable---here $x$, pulling the common factor out of the linear term.}
\end{codepairs}

The deeper operation is \B{Factor}, which writes a polynomial as a product of\index{polynomial!factorization}
irreducible pieces. What counts as irreducible\index{polynomial!irreducible} depends on the number system you allow,
and Mathilda lets you choose:

\begin{codepairs}
\pair{03-introduction/algebra-factor/1}{Over the rationals, $x^{6}-1$ splits into four
factors.}
\pair{03-introduction/algebra-factor/2}{But $x^{2}+1$ will not factor at all---it is
irreducible over~$\mathbb{Q}$.}
\pair{03-introduction/algebra-factor/3}{Allow the imaginary unit, with
\texttt{Extension}, and it becomes\index{field extension} $(x-i)(x+i)$.}
\pair{03-introduction/algebra-factor/4}{Allow $\sqrt2$ and $x^{4}+1$ splits into two
real quadratics. Extensions return in Section~\ref{sec:algebra}.}
\end{codepairs}

Closely related to factoring is the polynomial greatest common divisor\index{polynomial!greatest common divisor}. Just as \B{GCD}
finds the largest integer dividing two whole numbers, \B{PolynomialGCD} finds the
highest-degree polynomial dividing two polynomials---their shared factor:

\begin{codepairs}
\pair{03-introduction/algebra-gcd/1}{These two polynomials share the factor $x-1$.}
\pair{03-introduction/algebra-gcd/2}{These share $x^{2}-1$, found directly, without
factoring either polynomial.}
\end{codepairs}

That common factor is exactly what simplifies a \emph{rational function}\index{rational function}, a ratio of
polynomials. Mathilda provides the three everyday operations:

\begin{codepairs}
\pair{03-introduction/algebra-rational/1}{\B{Cancel} divides out the common factor
(the $x-1$ above), reducing the fraction to lowest terms.}
\pair{03-introduction/algebra-rational/2}{\B{Together} combines a sum of fractions
over one common denominator.}
\pair{03-introduction/algebra-rational/3}{\B{Apart} does the reverse---the partial-%
fraction decomposition you reach for\index{partial fraction decomposition} before integrating. More in
Section~\ref{sec:algebra}.}
\end{codepairs}

Beyond a single polynomial lies a whole \emph{system} of them, and the tool for those is
the \B{GroebnerBasis}: a canonical set of generators for the ideal\index{ideal} a system defines. It
is the multivariate, nonlinear analogue of Gaussian elimination\index{Gaussian elimination}---it rewrites a system
into an equivalent triangular one that is far easier to solve, and it can eliminate
chosen variables outright.

\begin{codepairs}
\pair{03-introduction/algebra-groebner/1}{The unit circle meets the parabola $y=x^{2}$.
\B{GroebnerBasis} triangulates the system: the first polynomial involves only $y$---a
quadratic to solve---and the second then recovers $x$ from $y$.}
\pair{03-introduction/algebra-groebner/2}{Naming a variable to eliminate---here the
parameter $t$ of the curve $(t^{2},t^{3})$---collapses the system to the curve's implicit
equation $x^{3}=y^{2}$, the cuspidal cubic. Section~\ref{sec:algebra} returns to
Gr\"obner bases.}
\end{codepairs}

\section{Solving equations}
\label{sec:intro-solve}

Solving equations is among the oldest tasks in mathematics, and Mathilda does it both
exactly and numerically.

\subsection{Exact solutions of equations}

\B{Solve} finds \emph{exact} solutions---as integers, radicals, or general closed
forms---for a single equation or a whole system:

\begin{codepairs}
\pair{03-introduction/solve-exact/1}{A quadratic with two integer roots. Note
\texttt{==}, which builds an equation\index{equation vs assignment}, versus \texttt{=}, which assigns.}
\pair{03-introduction/solve-exact/2}{$x^{2}=2$: the two exact surds $\pm\sqrt2$.}
\pair{03-introduction/solve-exact/3}{A linear system in two unknowns\index{system of equations!linear}, uniquely solved.}
\pair{03-introduction/solve-exact/4}{A \emph{nonlinear} system---a circle meeting the
line $y=x$---with both intersection points, exactly.}
\pair{03-introduction/solve-exact/5}{A cubic with no rational roots: the exact solutions
come back as \B{Root} objects, each an implicit algebraic number\index{algebraic number}---the $k$th root of the
given polynomial---rather than an unwieldy nest of radicals.}
\pair{03-introduction/solve-exact/6}{A \B{Root} object is still an exact number, so \B{N}
evaluates it to as many digits as you ask---here to machine precision.}
\end{codepairs}

\subsection{Numerical solutions of equations}

When an exact solution is unavailable---the general quintic has no formula in radicals,
by Abel's theorem\index{Abel's theorem}---or simply not wanted, the numerical solvers take over. \B{NRoots}
and \B{NSolve} return every root of a polynomial equation or system as machine numbers,
while \B{FindRoot} follows Newton's method\index{Newton's method} from a starting point and handles
transcendental problems too:

\begin{codepairs}
\pair{03-introduction/solve-numerical/1}{\B{NRoots} finds all five roots of $x^{5}-x-1$:
one real and two complex-conjugate pairs, none expressible in radicals.}
\pair{03-introduction/solve-numerical/2}{\B{NSolve} handles systems. Eliminating $y$ here
leaves that very same quintic, so the five solutions carry its unsolvable roots in their
$x$-coordinate.}
\pair{03-introduction/solve-numerical/3}{\B{FindRoot}, started near $x=1$, converges to
the root of $\cos x = x$, the ``Dottie number''\index{Dottie number}---a transcendental equation with no
closed form.}
\pair{03-introduction/solve-numerical/4}{\B{FindRoot} takes systems as well: a
transcendental pair in $x$ and $y$, solved from the starting point $(1,1)$.}
\end{codepairs}

\subsection{Equations with a continuum of solutions}

\B{Solve} expects a finite list of solutions to name. But an equation built from
absolute values or nested radicals can hold across a whole interval, and then what
is wanted is again a \emph{description} of the solution set\index{solution
set}\index{equation!continuum of solutions} rather than a list of points.
\B{Reduce}---the general real-equation engine taken up in the next
section---provides it, solving over the reals and returning the exact region on
which the equation holds:

\begin{codepairs}
\pair{03-introduction/solve-region/1}{Each nested radical\index{radical!nested} is a
perfect square in disguise---$x+1-2\sqrt x=(\sqrt x-1)^{2}$ and
$x+16-8\sqrt x=(\sqrt x-4)^{2}$---so the left side is $|\sqrt x-1|+|\sqrt x-4|$,
which equals~$3$ across the entire interval $1\le x\le16$.}
\pair{03-introduction/solve-region/2}{A sum of absolute values\index{absolute value}
over~$x$: the equation is satisfied not at isolated points but along the whole ray
$x\ge3$.}
\pair{03-introduction/solve-region/3}{Nested absolute values, even in~$x$: the
solution set is a \emph{union} of two intervals, one mirrored across the origin from
the other.}
\end{codepairs}

Exact solving belongs to the algebra of Section~\ref{sec:algebra}; the numerical
solvers to Section~\ref{sec:numcalc}.

\section{Solving inequalities}
\label{sec:intro-inequalities}

An equation has a handful of solutions to name; an inequality\index{inequality} usually
has infinitely many, so what you want back is not a list but a \emph{description} of the
solution set\index{solution set}. \B{Reduce} provides it: given one or more polynomial,
rational, or radical inequalities it returns an equivalent, exhaustive statement---an
interval, a union of intervals, or a region of the plane---computed by \emph{cylindrical
algebraic
decomposition}\index{cylindrical algebraic decomposition}, a complete decision
procedure\index{decision procedure} for systems of polynomial equations and inequalities
over the reals.

\begin{codepairs}
\pair{03-introduction/inequalities/1}{A quadratic inequality: the solution set is the
closed interval where the parabola lies on or below the axis.}
\pair{03-introduction/inequalities/2}{A quartic can cut the line into several
pieces---here a union of two open intervals.}
\pair{03-introduction/inequalities/3}{A rational inequality. The pole at $x=-2$ is
excluded automatically---\B{Reduce} keeps track of where the expression is undefined.}
\pair{03-introduction/inequalities/4}{The endpoints need not be rational: here they are
the exact surds $\pm\sqrt2$.}
\pair{03-introduction/inequalities/5}{A rational inequality whose numerator carries the
denominator's factors---$x^{3}-x=x(x^{2}-1)$---so after cancellation it collapses to the
plain linear condition $x\le2$.}
\pair{03-introduction/inequalities/6}{A \emph{radical}
inequality\index{inequality!radical}, not polynomial at all. The square root's domain
$x\ge-4$ meets the single crossing at $x=5$, where $\sqrt{x+4}=x-2$, to give
$-4\le x<5$.}
\end{codepairs}

Because it \emph{decides} the theory rather than merely searching for solutions,
\B{Reduce} settles the universal questions too\calloutnote{Under the hood}{That the real
numbers admit \emph{any} such decision procedure is Tarski's theorem; cylindrical
algebraic decomposition is the algorithm that makes it practical, and it is what lets
\B{Reduce} answer a question about \emph{all} real $x$ with a definite \texttt{True} or
\texttt{False}.}, returning \texttt{True} or \texttt{False} outright when the solution
set turns out to be everything, or nothing:

\begin{codepairs}
\pair{03-introduction/inequalities/7}{$x^{2}+1>0$ holds for every real $x$---a one-line
proof that the expression is positive-definite.}
\pair{03-introduction/inequalities/8}{No real number makes $x^{2}+x+1$ negative (its
discriminant is negative), so the answer is \texttt{False}.}
\pair{03-introduction/inequalities/9}{The arithmetic--geometric-mean
inequality\index{AM--GM inequality} $x+1/x\ge2$: \B{Reduce} confirms it across the
\emph{whole} positive ray, the constraint collapsing to just $x>0$.}
\pair{03-introduction/inequalities/10}{With two variables the solution set becomes a
region of the plane---the open unit disk, described as $x$ ranging over $(-1,1)$ with $y$
between the two branches of its bounding circle.}
\pair{03-introduction/inequalities/11}{Absolute values\index{absolute value} in two
variables: the open \emph{diamond} $|x|+|y|<1$, the unit ball of the $\ell^{1}$
norm\index{norm!$\ell^1$}. \B{Reduce} case-splits each \B{Abs} on the sign of its
argument, so the region falls out as its four quadrant sectors.}
\end{codepairs}

\section{Simplification}
\label{sec:intro-simplify}

Expanding, factoring, and cancelling each drive an expression toward a \emph{particular}
form. Often, though, you simply want the \emph{simplest} form, whatever that turns out
to be. \B{Simplify} searches for it\index{simplification}, trying a battery of algebraic and trigonometric
transformations and keeping the smallest result:

\begin{codepairs}
\pair{03-introduction/simplification/1}{The Pythagorean identity\index{Pythagorean identity} collapses to~$1$.}
\pair{03-introduction/simplification/2}{A rational function reduces.}
\pair{03-introduction/simplification/3}{A double-angle identity is recognised:
$\cos^{2}x-\sin^{2}x=\cos 2x$.}
\pair{03-introduction/simplification/4}{And a nested radical
\emph{denests}\index{radical!denesting}: $\sqrt{3+2\sqrt2}=1+\sqrt2$.}
\end{codepairs}

Some simplifications hold only under conditions, and Mathilda will not assume what you
have not told it. Because $\sqrt{x^{2}}$ equals $|x|$, not $x$:

\begin{codepairs}
\pair{03-introduction/simplification-assumptions/1}{With nothing known about $x$,
$\sqrt{x^{2}}$ is left exactly as written.}
\pair{03-introduction/simplification-assumptions/2}{Told that $x>0$, it simplifies to
$x$.}
\end{codepairs}

An assumption can be more interesting\index{assumptions} than a sign condition---sometimes what matters is
the \emph{kind} of number involved:

\begin{codepairs}
\pair{03-introduction/simplification-integer/1}{For a general~$n$, $\sin(n\pi)$ is left
untouched.}
\pair{03-introduction/simplification-integer/2}{Told that $n$ is an integer, with
\B{Element}, Mathilda proves it must vanish. More in Section~\ref{sec:algebra}.}
\end{codepairs}

\section{Calculus}
\label{sec:intro-calculus}

Mathilda does symbolic calculus\index{symbolic calculus}: differentiation, limits, and integration on
expressions, returning exact results. Differentiation with \B{D} applies the rules of
calculus recursively, and by naming a second variable it takes partial derivatives\index{partial derivative}:

\begin{codepairs}
\pair{03-introduction/calculus-diff/1}{The power rule, applied term by term.}
\pair{03-introduction/calculus-diff/2}{The product and chain rules, automatically.}
\pair{03-introduction/calculus-diff/3}{Symbolic exponents are fine too.}
\pair{03-introduction/calculus-diff/4}{A \emph{partial} derivative in $x$, treating
$y$ as a constant.}
\pair{03-introduction/calculus-diff/5}{The same expression, now differentiated in
$y$.}
\end{codepairs}

Limits are computed with\index{limit} \B{Limit}, including the classic ones that define the
derivative and the number~$e$:

\begin{codepairs}
\pair{03-introduction/calculus-limits/1}{The limit at the heart of the derivative of
sine.}
\pair{03-introduction/calculus-limits/2}{One of the definitions of~$e$.}
\pair{03-introduction/calculus-limits/3}{The leading behaviour of $1-\cos x$ near the
origin.}
\end{codepairs}

Near a point a function is captured by its power series\index{power series}, and \B{Series} computes it to any
order, returning the polynomial part together with an $O[x]^{n}$ term that records where
the expansion was truncated:

\begin{codepairs}
\pair{03-introduction/calculus-series/1}{The exponential series, to sixth order.}
\pair{03-introduction/calculus-series/2}{$\arctan$: the alternating odd reciprocals whose
value at $x=1$ gives Leibniz's series for $\pi/4$.}
\pair{03-introduction/calculus-series/3}{At a pole the expansion is a \emph{Laurent}\index{Laurent series}
series---here $\cot x$ begins with a $1/x$ term.}
\pair{03-introduction/calculus-series/4}{With a \emph{symbolic} exponent the coefficients
come out as the generalized binomial coefficients.}
\pair{03-introduction/calculus-series/5}{\B{Normal} drops the $O[\,]$ term, turning the
\texttt{SeriesData} object\index{SeriesData} a series is into an ordinary polynomial.}
\end{codepairs}

Integration, the harder inverse of differentiation\index{integration}, is done with \B{Integrate}---%
indefinite integrals as symbolic antiderivatives, and definite integrals\index{definite integral}, including
improper ones:

\begin{codepairs}
\pair{03-introduction/calculus-integrate/1}{An indefinite integral by parts.}
\pair{03-introduction/calculus-integrate/2}{Another, giving an inverse tangent.}
\pair{03-introduction/calculus-integrate/3}{A definite integral over $[0,1]$.}
\pair{03-introduction/calculus-integrate/4}{The Gaussian integral over the whole real\index{Gaussian integral}
line: $\int_{-\infty}^{\infty} e^{-x^{2}}\,dx=\sqrt\pi$.}
\end{codepairs}

\B{Integrate} is not confined to one variable. Give it several iterators and it
computes a \emph{multiple} integral\index{multiple integral}, from the inside out, so the inner limits may
depend on an outer variable:

\begin{codepairs}
\pair{03-introduction/calculus-integrate-multi/1}{$x y$ integrated over the unit
square.}
\pair{03-introduction/calculus-integrate-multi/2}{Over the triangle $0\le y\le x\le 1$,
whose inner bound is $x$ itself.}
\end{codepairs}

And when the limits are \emph{complex}, \B{Integrate} evaluates a \emph{contour\index{contour integral}
integral} along the straight segments joining the points you list:

\begin{codepairs}
\pair{03-introduction/calculus-integrate-contour/1}{A line integral of $z^{2}$ along
the segment from $0$ to $1+i$.}
\pair{03-introduction/calculus-integrate-contour/2}{A closed loop around the origin:
$\oint \frac{dz}{z}=2\pi i$, the cornerstone of the residue calculus\index{residue calculus}.
Section~\ref{sec:calculus} has the full story.}
\end{codepairs}

\section{Numerical calculus}
\label{sec:intro-numcalc}

Not every problem has a closed form, and even when it does a number is sometimes all
you want. The \texttt{N}-prefixed family answers\index{numerical calculus} these questions numerically and
adaptively---to machine precision by default, or to any arbitrary precision you request
with the \texttt{WorkingPrecision} option\index{WorkingPrecision option}:

\begin{codepairs}
\pair{03-introduction/numerical-calculus/1}{\B{ND} takes a numerical derivative---here
of $\sin$ at $1$, giving $\cos 1$.}
\pair{03-introduction/numerical-calculus/2}{\B{NIntegrate} evaluates the Gaussian
integral numerically ($\approx\sqrt\pi$).}
\pair{03-introduction/numerical-calculus/3}{\B{NSum} sums a series numerically.}
\end{codepairs}

Differential equations are solved numerically\index{differential equation} by \B{NDSolve}. Take a damped harmonic
oscillator, $y'' + \tfrac{3}{10}\,y' + 4y = 0$ with $y(0)=1$ and $y'(0)=0$: the solution
comes back as an \B{InterpolatingFunction}, and \B{Plot} takes that
directly,\footnote{That plot is genuinely the \texttt{Out[2]} of the \B{Plot} command,
drawn by the running system. To keep it, \B{Export} writes any such graphic to PDF, PNG,
or JPEG (\texttt{Export["decay.pdf", plot]})---indeed every figure in this book is
produced exactly this way.} so the motion---a decaying oscillation settling to rest---is
one further command away:

\plotcell{03-introduction/ndsolve-solution}

Finally, \B{NLimit} evaluates a limit numerically, invaluable when a symbolic limit is
hard to obtain:

\begin{codepairs}
\pair{03-introduction/numerical-calculus-nlimit/1}{The familiar limit that
defines~$e$.}
\pair{03-introduction/numerical-calculus-nlimit/2}{More striking: $\zeta(s)$ (the
\B{Zeta} function) and $1/(s-1)$ both diverge at $s=1$, yet their difference tends to
the Euler--Mascheroni constant $\gamma$.}
\pair{03-introduction/numerical-calculus-nlimit/3}{A slower, non-obvious limit,
recovered numerically.}
\end{codepairs}

Section~\ref{sec:numcalc} covers numerical differentiation, integration, summation,
differential equations, and limits.

\section{Linear algebra}
\label{sec:intro-linalg}

Vectors and matrices are not a new kind of object in Mathilda---a
vector\index{vector} is just a list, and a matrix\index{matrix} a list of its
rows---so everything you already know about lists carries over unchanged. What
linear algebra adds is a vocabulary of operations on them, beginning with the
\emph{dot product}\index{dot product}, written with an ordinary period. The single
operator \texttt{.}\ (\B{Dot}) serves throughout: it contracts two vectors to a
scalar, applies a matrix to a vector, and multiplies two matrices.

\begin{codepairs}
\pair{03-introduction/linear-algebra-vectors/1}{The dot product of two vectors is a
single number.}
\pair{03-introduction/linear-algebra-vectors/2}{\B{Norm} is the Euclidean
length\index{norm!Euclidean} $\sqrt{v\cdot v}$, kept exact---here the surd
$\sqrt{14}$.}
\pair{03-introduction/linear-algebra-vectors/3}{\B{Cross} is the three-dimensional
cross product\index{cross product}: $\hat\imath\times\hat\jmath=\hat k$.}
\pair{03-introduction/linear-algebra-vectors/4}{The same \texttt{.}\ applies a matrix
to a vector\index{matrix!times vector}.}
\pair{03-introduction/linear-algebra-vectors/5}{And multiplies two
matrices\index{matrix!product}---here by the swap matrix, which exchanges the
columns.}
\end{codepairs}

The staple operations on a square matrix are each a single call away, and---this is
the point---they return \emph{exact} answers. Working over the rationals rather than
in floating point, \B{Inverse} produces an exact rational matrix, and multiplying it
back recovers the identity with no rounding error at all:

\begin{codepairs}
\pair{03-introduction/linear-algebra-matrices/1}{A matrix is a list of its rows; \texttt{A}
is defined once and reused below.}
\pair{03-introduction/linear-algebra-matrices/2}{\B{Det} is the
determinant\index{determinant}.}
\pair{03-introduction/linear-algebra-matrices/3}{\B{Tr} is the trace\index{trace}, the
sum of the diagonal entries.}
\pair{03-introduction/linear-algebra-matrices/4}{\B{Inverse} gives the exact inverse,
with rational entries.}
\pair{03-introduction/linear-algebra-matrices/5}{Multiplying a matrix by its inverse
returns the identity\index{identity matrix}, exactly.}
\pair{03-introduction/linear-algebra-matrices/6}{\B{Transpose} reflects a matrix across
its diagonal, turning rows into columns.}
\pair{03-introduction/linear-algebra-matrices/7}{\B{IdentityMatrix} builds the
$n\times n$ identity.}
\pair{03-introduction/linear-algebra-matrices/8}{\B{DiagonalMatrix} places a list along
the diagonal, zeros elsewhere.}
\end{codepairs}

Nothing here requires the entries to be numbers. Given symbols, the same functions
return the general formula\index{matrix!symbolic}---the textbook determinant and inverse
of a $2\times2$ matrix:

\begin{codepairs}
\pair{03-introduction/linear-algebra-symbolic/1}{The determinant $ad-bc$.}
\pair{03-introduction/linear-algebra-symbolic/2}{And the inverse, each entry divided by
that determinant.}
\end{codepairs}

\subsection{Linear systems}

Solving $A\mathbf{x}=\mathbf{b}$ is the oldest task in the subject. \B{LinearSolve}
does it directly; \B{RowReduce} performs the Gaussian elimination\index{Gaussian
elimination} behind it, returning the reduced row echelon form\index{row echelon form};
and \B{MatrixRank} and \B{NullSpace} report the structure of the solution set---the
rank\index{rank (of a matrix)} of the coefficient matrix and a basis for the vectors it
sends to zero.

\begin{codepairs}
\pair{03-introduction/linear-algebra-solve/1}{\B{LinearSolve} solves a $2\times2$ system
exactly.}
\pair{03-introduction/linear-algebra-solve/2}{\B{RowReduce} drives a matrix to reduced row
echelon form; this one is invertible, so the result is the identity.}
\pair{03-introduction/linear-algebra-solve/3}{\B{MatrixRank}: these two rows are
proportional, so the rank is~$1$.}
\pair{03-introduction/linear-algebra-solve/4}{\B{NullSpace} returns a basis for the
vectors the matrix annihilates.}
\end{codepairs}

\subsection{Eigenvalues and eigenvectors}

The deepest invariants of a matrix are its \emph{eigenvalues}\index{eigenvalue} and
\emph{eigenvectors}\index{eigenvector}---the scalars $\lambda$ and directions
$\mathbf{v}$ for which $A\mathbf{v}=\lambda\mathbf{v}$. \B{Eigenvalues} and
\B{Eigenvectors} compute them exactly:

\begin{codepairs}
\pair{03-introduction/linear-algebra-eigen/1}{A symmetric matrix with integer
eigenvalues.}
\pair{03-introduction/linear-algebra-eigen/2}{Its eigenvectors, orthogonal as symmetry
demands.}
\pair{03-introduction/linear-algebra-eigen/3}{Here the eigenvalues are the golden
ratio\index{golden ratio} $\varphi=(1+\sqrt5)/2$ and its conjugate---the growth rates
concealed in the Fibonacci recurrence.}
\pair{03-introduction/linear-algebra-eigen/4}{Indeed the powers of that same matrix
\emph{are} the Fibonacci numbers\index{Fibonacci numbers}: the tenth power displays
$F_{11}=89$ and $F_{10}=55$.}
\pair{03-introduction/linear-algebra-eigen/5}{When the eigenvalues have no radical form
they come back as \B{Root} objects\index{Root object}---the roots of $x^{3}-x-1$, exact
algebraic numbers, exactly as \B{Solve} returned a cubic in \S\ref{sec:intro-solve}.}
\pair{03-introduction/linear-algebra-eigen/6}{Numericalising those exact eigenvalues
with \texttt{N} resolves the real root to $\approx1.32472$ and the complex-conjugate
pair to $-0.662359\pm0.56228\,i$.}
\pair{03-introduction/linear-algebra-eigen/7}{Asking instead for the eigenvalues of the
floating-point matrix lands on the identical list---by a numerical algorithm that never
forms the characteristic polynomial at all.}
\end{codepairs}

Those last two lines reach the \emph{same} three numbers by opposite roads. The first
takes the exact \Bnoidx{Root} eigenvalues and rounds them to machine precision; the second
hands the floating-point matrix straight to a numerical eigenvalue algorithm, which
never builds a characteristic polynomial\index{characteristic polynomial} to
factor---though \B{CharacteristicPolynomial} forms that polynomial, $\det(m-xI)$
(or $\det(m-xa)$ for a pencil $\{m,a\}$), on demand. That the two agree is a reassuring
cross-check---and a first sight of a choice that runs through all of computational
linear algebra\index{linear algebra!exact versus numerical}: exact or numerical, and
what each costs.

Section~\ref{sec:linalg} takes up linear algebra---exact and numerical alike---in
depth.

\section{Programming}
\label{sec:intro-programming}

Mathilda is also a programming language\index{programming language}, and an unusually expressive one. It supports
three styles that blend freely: pattern-based rewriting, procedural code, and
functional programming.

\subsection*{Pattern matching}

The most characteristic style is definition by \emph{patterns}\index{pattern matching}. A name followed by an
underscore, like \texttt{x\_}, matches anything and binds it\index{pattern matching!Blank (underscore)}, and a pattern can be
narrowed: \texttt{x\_Integer} matches only an expression with head \texttt{Integer}\index{head (of expression)},
and \texttt{?Positive} adds a test\index{pattern matching!test} the match must pass. A definition with \texttt{:=}
applies whenever its pattern fits:

\begin{codepairs}
\pair{03-introduction/pattern-matching/1}{Define \texttt{f} only on positive
integers.}
\pair{03-introduction/pattern-matching/2}{A positive integer matches: $7^{2}=49$.}
\pair{03-introduction/pattern-matching/3}{A negative integer does not match.}
\pair{03-introduction/pattern-matching/4}{A real does not match either.}
\pair{03-introduction/pattern-matching/5}{Nor does a complex number; each call is
returned unevaluated.}
\pair{03-introduction/pattern-matching/6}{\B{Cases} keeps the elements matching a
pattern---the even ones.}
\pair{03-introduction/pattern-matching/7}{The rule \texttt{/.}\ (\B{ReplaceAll})
substitutes into an expression.}
\pair{03-introduction/pattern-matching/8}{A pattern can destructure a two-element
list.}
\pair{03-introduction/pattern-matching/9}{It binds both parts at once. More in
Chapter~\ref{ch:programming}.}
\end{codepairs}

\subsection*{Procedural programming}

When an ordinary loop is the clearest thing, Mathilda has loops. \B{Module} introduces
local variables\index{local variable (Module)}, and \B{Do}, \B{For}, and \B{While} iterate:

\begin{codepairs}
\pair{03-introduction/procedural/1}{A \B{Do} loop sums $1$ through $100$.}
\pair{03-introduction/procedural/2}{A \B{For} loop computes $5! = 120$.}
\pair{03-introduction/procedural/3}{A \B{While} loop counts the steps of the Collatz
iteration from $27$.}
\end{codepairs}

\subsection*{Functional programming}

Often the shortest code is functional\index{functional programming}---building a result by applying and composing
functions rather than by mutating variables. The \texttt{\&} notation makes an
anonymous function, with \texttt{\#} as its argument\index{anonymous function}:

\begin{codepairs}
\pair{03-introduction/functional/1}{\B{Map} applies a function to every element.}
\pair{03-introduction/functional/2}{\B{Select} keeps those passing a test---the
primes.}
\pair{03-introduction/functional/3}{\B{Fold} chains a function across a list.}
\pair{03-introduction/functional/4}{\B{Nest} applies one repeatedly: $\sqrt{\cdot}$ of
$2$, three times.}
\pair{03-introduction/functional/5}{\B{Table} builds a list from a formula.}
\pair{03-introduction/functional/6}{\B{Apply}, written \texttt{@@}, replaces the head:
\texttt{Times @@ list} is the product.}
\pair{03-introduction/functional/7}{\texttt{@@@} applies at level one, to each
sublist.}
\pair{03-introduction/functional/8}{\B{Cases} picks out the elements matching a
pattern---the integers.}
\end{codepairs}

\section{Graphics}
\label{sec:intro-graphics}

Mathilda draws its own graphics. Every plotting function returns a \texttt{Graphics}\index{Graphics object} (in two
dimensions) or \texttt{Graphics3D} (in three) object built from primitives---lines, points,
polygons, colours---which the built-in engine renders on screen and \B{Export} writes to
PDF, PNG, or JPEG. This closing tour gives one example of each plotting function; as in
\S\ref{sec:intro-numcalc}, every plot below is the genuine \texttt{Out} of the command
shown, produced by the running system.

\subsection{Functions and data}

\B{Plot} draws one or more functions of a real variable---here the first odd harmonics of
a square wave:

\plotcell{03-introduction/graphics/plot}

\noindent \B{ListPlot} draws discrete data as points:

\plotcell{03-introduction/graphics/listplot}

\noindent \B{ParametricPlot} traces a curve $(x(t), y(t))$---a Lissajous figure---and
\B{PolarPlot} draws $r(\theta)$, here a five-petal rose:

\plotcell{03-introduction/graphics/parametricplot}

\plotcell{03-introduction/graphics/polarplot}

\subsection{Vector fields}

\B{VectorPlot} shows a field as a grid of arrows\index{vector field}, \B{StreamPlot} as flow lines. Here is the
phase field of a pendulum, then a pure rotation:

\plotcell{03-introduction/graphics/vectorplot}

\plotcell{03-introduction/graphics/streamplot}

\subsection{Scalar fields and arrays}

\B{ContourPlot} draws level curves of $f(x,y)$, \B{DensityPlot} shades the plane by value,
and \B{ArrayPlot} renders a matrix as a grid of cells---here the greatest-common-divisor
table, whose diagonals of bright cells expose its number-theoretic structure:

\plotcell{03-introduction/graphics/contourplot}

\plotcell{03-introduction/graphics/densityplot}

\plotcell{03-introduction/graphics/arrayplot}

\subsection{Charts and graphs}

\B{BarChart} draws labelled bars; \B{GraphPlot} lays out a graph automatically:

\plotcell{03-introduction/graphics/barchart}

\plotcell{03-introduction/graphics/graphplot}

\subsection{Complex functions}

\B{ComplexPlot} colours the plane by the value of a complex function\index{domain coloring}---hue for the phase,
brightness for the magnitude---so its zeros and poles stand out at a glance:

\plotcell{03-introduction/graphics/complexplot}

\subsection{Three dimensions}

\B{Plot3D} draws a surface $z=f(x,y)$, \B{ParametricPlot3D} a curve or surface in space,
and \B{ComplexPlot3D} the modulus of a complex function as a surface coloured by phase.
These return \texttt{Graphics3D}, which the engine renders with lighting and a movable camera,
and which \B{Export} writes as PNG or JPEG:

\plotcell{03-introduction/graphics/plot3d}

\plotcell{03-introduction/graphics/parametricplot3d}

\plotcell{03-introduction/graphics/complexplot3d}

\bigskip
\noindent
That is the tour. Every area glimpsed here has a chapter of its own ahead, and this
introduction will itself grow as Mathilda does. With the shape of the system in view,
we can now begin in earnest---starting, in the next part, with numbers.
